\documentclass[12pt,a4paper]{article}

% --- Metadata ---
\def\courseshortheader{HỌC MÁY ỨNG DỤNG --- NUMPY}
\def\coverbookline{TRÍ TUỆ NHÂN TẠO}
\def\covermaintitle{NumPy UltraQuick Tutorial}
\def\coversubtitle{Dịch tiếng Việt từ Google ML Crash Course (Colab)}

% Shared preamble for nhom_5 (requires \courseshortheader before \input)
\usepackage[utf8]{inputenc}
\usepackage[vietnamese]{babel}
\usepackage{amsmath,amssymb,amsthm}
\usepackage{geometry}
\usepackage{enumitem}
\usepackage[most]{tcolorbox}
\usepackage{hyperref}
\usepackage{fancyhdr}
\usepackage{graphicx}
\usepackage{booktabs}
\usepackage{tikz}
\usetikzlibrary{calc,arrows.meta,decorations.pathreplacing,positioning}
\usepackage{eso-pic}
\usepackage{xcolor}

\geometry{a4paper, margin=2cm, bottom=2.5cm}
\hypersetup{colorlinks=true, linkcolor=blue!70!black, urlcolor=blue}
\setlist[itemize]{topsep=4pt, itemsep=2pt}
\setlist[enumerate]{topsep=4pt, itemsep=2pt}

\AddToShipoutPictureFG{%
  \AtPageCenter{%
    \begin{tikzpicture}[remember picture, overlay]
      \node[rotate=45, scale=1, opacity=0.06]{%
        \fontsize{60}{72}\selectfont\bfseries\sffamily Đặng Tấn Quí%
      };
    \end{tikzpicture}%
  }%
}

\pagestyle{fancy}
\fancyhf{}
\fancyhead[L]{\small\textit{\courseshortheader}}
\fancyhead[R]{\small\textit{Đặng Tấn Quí}}
\fancyfoot[C]{\thepage}
\fancyfoot[L]{\footnotesize\textcopyright\ Đặng Tấn Quí}
\fancyfoot[R]{\footnotesize SĐT: 0377\,122\,210}
\renewcommand{\headrulewidth}{0.4pt}
\renewcommand{\footrulewidth}{0.4pt}

\fancypagestyle{plain}{%
  \fancyhf{}
  \fancyfoot[C]{\thepage}
  \fancyfoot[L]{\footnotesize\textcopyright\ Đặng Tấn Quí}
  \fancyfoot[R]{\footnotesize SĐT: 0377\,122\,210}
  \renewcommand{\headrulewidth}{0pt}
  \renewcommand{\footrulewidth}{0.4pt}
}

\newtcolorbox{dongco}[1][]{%
  enhanced, breakable,
  colback=teal!4!white, colframe=teal!60!black,
  fonttitle=\bfseries, title={Động cơ học -- Tại sao cần khái niệm này?}, #1%
}
\newtcolorbox{ynghia}[1][]{%
  enhanced, breakable,
  colback=purple!3!white, colframe=purple!50!black,
  fonttitle=\bfseries, title={Ý nghĩa \& Trực giác}, #1%
}
\newtcolorbox{ungdung}[1][]{%
  enhanced, breakable,
  colback=olive!5!white, colframe=olive!60!black,
  fonttitle=\bfseries, title={Ứng dụng thực tế}, #1%
}
\newtcolorbox{dinhnghia}[1][]{%
  enhanced, breakable,
  colback=blue!3!white, colframe=blue!60!black,
  fonttitle=\bfseries, title={Định nghĩa}, #1%
}
\newtcolorbox{dinhly}[1][]{%
  enhanced, breakable,
  colback=green!3!white, colframe=green!50!black,
  fonttitle=\bfseries, title={Định lý}, #1%
}
\newtcolorbox{tinhchat}[1][]{%
  enhanced, breakable,
  colback=yellow!5!white, colframe=orange!50!black,
  fonttitle=\bfseries, title={Tính chất}, #1%
}
\newtcolorbox{phuongphap}[1][]{%
  enhanced, breakable,
  colback=gray!5!white, colframe=gray!50!black,
  fonttitle=\bfseries, title={Phương pháp}, #1%
}
\newtcolorbox{luuy}[1][]{%
  enhanced, breakable,
  colback=red!3!white, colframe=red!50!black,
  fonttitle=\bfseries, title={Lưu ý}, #1%
}
\newtcolorbox{congthuc}[1][]{%
  enhanced, breakable,
  colback=cyan!3!white, colframe=cyan!50!black,
  fonttitle=\bfseries, title={Công thức}, #1%
}
\newtcolorbox{vidu}[1][]{%
  enhanced, breakable,
  colback=violet!3!white, colframe=violet!50!black,
  fonttitle=\bfseries, title={Ví dụ}, #1%
}
\newtcolorbox{phanvidu}[1][]{%
  enhanced, breakable,
  colback=red!2!white, colframe=red!40!black,
  fonttitle=\bfseries, title={Phản ví dụ}, #1%
}
\newtcolorbox{tongket}[1][]{%
  enhanced, breakable,
  colback=blue!4!white, colframe=blue!70!black,
  fonttitle=\bfseries, title={Tổng kết chương}, #1%
}


\usepackage{listings}
\usepackage{upquote}
\usepackage[T1]{fontenc}

\lstdefinestyle{mlcode}{
  basicstyle=\ttfamily\small,
  breaklines=true,
  breakatwhitespace=false,
  columns=fullflexible,
  keepspaces=true,
  showstringspaces=false,
  frame=single,
  framerule=0.4pt,
  rulecolor=\color{black!25},
  backgroundcolor=\color{black!2},
  xleftmargin=0.2cm,
  xrightmargin=0.2cm
}
\lstset{style=mlcode}

\newcommand{\mlccurl}[1]{\url{#1}}

\begin{document}

\thispagestyle{empty}
\begin{tikzpicture}[remember picture, overlay]
  \fill[blue!80!black] (current page.south west) rectangle (current page.north east);
  \fill[blue!60!cyan, opacity=0.8]
    (current page.south west) -- (current page.north west) --
    ([xshift=-3cm]current page.north east) -- (current page.south east) -- cycle;
  \foreach \r/\o in {6/0.05, 4.5/0.08, 3/0.12} {
    \fill[white, opacity=\o] ([xshift=2cm, yshift=-2cm]current page.north east) circle (\r cm);
  }
  \draw[white, line width=2pt] ([yshift=-5cm]current page.north west) ++(2cm,0) -- ++(17cm,0);
  \draw[white, line width=2pt] ([yshift=-16cm]current page.north west) ++(2cm,0) -- ++(17cm,0);
  \node[anchor=west, white, font=\fontsize{28}{34}\bfseries\selectfont]
    at ([xshift=3cm, yshift=-7.5cm]current page.north west) {\coverbookline};
  \node[anchor=west, white, font=\fontsize{20}{26}\selectfont]
    at ([xshift=3cm, yshift=-9cm]current page.north west) {\covermaintitle};
  \node[anchor=west, white, font=\fontsize{16}{22}\selectfont]
    at ([xshift=3cm, yshift=-10.2cm]current page.north west) {\coversubtitle};
  \node[anchor=west, white, font=\Large\bfseries]
    at ([xshift=3cm, yshift=-13cm]current page.north west) {Biên soạn:};
  \node[anchor=west, yellow!80!white, font=\fontsize{24}{30}\bfseries\selectfont]
    at ([xshift=3cm, yshift=-14.5cm]current page.north west) {ĐẶNG TẤN QUÍ};
  \node[anchor=west, white!90, font=\large]
    at ([xshift=3cm, yshift=-18cm]current page.north west) {SĐT: 0377\,122\,210};
  \node[anchor=south, white!80, font=\small]
    at ([yshift=1.5cm]current page.south) {Tài liệu có bản quyền --- Cấm sao chép khi chưa được sự đồng ý của tác giả};
  \node[anchor=south east, white, font=\Large\bfseries]
    at ([xshift=-2cm, yshift=3cm]current page.south east) {\the\year};
\end{tikzpicture}

\newpage
\tableofcontents
\newpage

%==============================================================================
\section*{Lời nói đầu}
\addcontentsline{toc}{section}{Lời nói đầu}
%==============================================================================

\begin{ynghia}[title={Nguồn}]
Tài liệu dịch từ notebook \textit{NumPy UltraQuick Tutorial} của Google Machine Learning Crash Course:
\mlccurl{https://colab.research.google.com/github/google/eng-edu/blob/main/ml/cc/exercises/numpy_ultraquick_tutorial.ipynb}
\end{ynghia}

\begin{luuy}[title={target}]
Đây không phải tutorial NumPy đầy đủ. target là cung cấp \textbf{vừa đủ} NumPy để làm các bài tập Colab trong ML Crash Course.
\end{luuy}

\newpage

%==============================================================================
\section{Colab là gì?}
%==============================================================================

\begin{dinhnghia}[title={Colaboratory (Colab)}]
Machine Learning Crash Course dùng \textbf{Colaboratory (Colab)} cho tất cả bài tập lập trình.
Colab là một hiện thực của Google dựa trên \textbf{Jupyter Notebook}:
\mlccurl{https://jupyter.org/}.

Để tìm hiểu cách dùng Colab, xem \textit{Welcome to Colaboratory}:
\mlccurl{https://research.google.com/colaboratory}.
\end{dinhnghia}

%==============================================================================
\section{Giới thiệu NumPy}
%==============================================================================

\begin{dongco}[title={Vì sao cần NumPy?}]
\textbf{NumPy} là thư viện Python để tạo và thao tác \textbf{ma trận} (cấu trúc dữ liệu trung tâm của nhiều thuật toán ML).
Ma trận là đối tượng toán học lưu giá trị theo \textbf{hàng} và \textbf{cột}.
\end{dongco}

\begin{ynghia}[title={Cách gọi ``ma trận'' trong các thư viện}]
Trong Python, ``ma trận'' thường được biểu diễn bằng \textbf{list}.
Trong NumPy, ``ma trận'' thường được gọi là \textbf{array}.
Trong TensorFlow, ``ma trận/mảng đa chiều'' thường được gọi là \textbf{tensor}.
\end{ynghia}

%==============================================================================
\section{Nhập (import) NumPy}
%==============================================================================

\begin{phuongphap}[title={Import module NumPy}]
Chạy cell sau để import NumPy.
\end{phuongphap}

\begin{lstlisting}
import numpy as np
\end{lstlisting}

%==============================================================================
\section{Khởi tạo array với giá trị cụ thể}
%==============================================================================

\begin{phuongphap}[title={Tạo array 1 chiều bằng \texttt{np.array}}]
Gọi \texttt{np.array} để tạo một NumPy array với các giá trị bạn chọn.
Ví dụ sau tạo một array 8 phần tử.
\end{phuongphap}

\begin{lstlisting}
one_dimensional_array = np.array([1.2, 2.4, 3.5, 4.7, 6.1, 7.2, 8.3, 9.5])
print(one_dimensional_array)
\end{lstlisting}

\begin{phuongphap}[title={Tạo array 2 chiều (ma trận)}]
Để tạo array 2 chiều, bạn đặt thêm một lớp ngoặc vuông. Ví dụ sau tạo ma trận $3 \times 2$.
\end{phuongphap}

\begin{lstlisting}
two_dimensional_array = np.array([[6, 5], [11, 7], [4, 8]])
print(two_dimensional_array)
\end{lstlisting}

\begin{luuy}[title={Array toàn số 0 và toàn số 1}]
Để tạo array toàn số 0, dùng \texttt{np.zeros}. Để tạo array toàn số 1, dùng \texttt{np.ones}.
\end{luuy}

%==============================================================================
\section{Khởi tạo array theo dãy số}
%==============================================================================

\begin{phuongphap}[title={Tạo dãy số bằng \texttt{np.arange}}]
Bạn có thể tạo array từ một dãy số liên tiếp.
\end{phuongphap}

\begin{lstlisting}
sequence_of_integers = np.arange(5, 12)
print(sequence_of_integers)
\end{lstlisting}

\begin{luuy}[title={Cận dưới và cận trên trong \texttt{np.arange}}]
\texttt{np.arange(5, 12)} tạo dãy \textbf{có} 5 nhưng \textbf{không} gồm 12.
\end{luuy}

%==============================================================================
\section{Khởi tạo array ngẫu nhiên}
%==============================================================================

\begin{phuongphap}[title={Số nguyên ngẫu nhiên bằng \texttt{np.random.randint}}]
\texttt{np.random.randint} sinh số nguyên ngẫu nhiên trong khoảng \([\texttt{low}, \texttt{high})\).
Ví dụ sau tạo array 6 phần tử trong khoảng 50 đến 100.
\end{phuongphap}

\begin{lstlisting}
random_integers_between_50_and_100 = np.random.randint(low=50, high=101, size=(6,))
print(random_integers_between_50_and_100)
\end{lstlisting}

\begin{luuy}[title={Giá trị lớn nhất sinh ra}]
Số lớn nhất mà \texttt{np.random.randint} có thể sinh ra luôn \textbf{nhỏ hơn} đối số \texttt{high}.
Vì vậy để có thể sinh 100, ta đặt \texttt{high=101}.
\end{luuy}

\begin{phuongphap}[title={Số thực ngẫu nhiên bằng \texttt{np.random.random}}]
\texttt{np.random.random} sinh số thực ngẫu nhiên trong \([0.0, 1.0)\).
\end{phuongphap}

\begin{lstlisting}
random_floats_between_0_and_1 = np.random.random((6,))
print(random_floats_between_0_and_1)
\end{lstlisting}

%==============================================================================
\section{Phép toán trên NumPy array và broadcasting}
%==============================================================================

\begin{dongco}[title={Vì sao cần broadcasting?}]
Trong đại số tuyến tính, phép cộng/trừ hai array thường yêu cầu cùng kích thước.
Phép nhân cũng có quy tắc tương thích kích thước chặt chẽ.
NumPy hỗ trợ một cơ chế gọi là \textbf{broadcasting} để ``mở rộng ảo'' toán hạng nhỏ hơn cho phù hợp kích thước.
\end{dongco}

\begin{phuongphap}[title={Cộng một hằng số (broadcasting)}]
Ví dụ sau cộng 2.0 vào mọi phần tử của array.
\end{phuongphap}

\begin{lstlisting}
random_floats_between_2_and_3 = random_floats_between_0_and_1 + 2.0
print(random_floats_between_2_and_3)
\end{lstlisting}

\begin{phuongphap}[title={Nhân với một hằng số (broadcasting)}]
Ví dụ sau nhân mọi phần tử với 3.
\end{phuongphap}

\begin{lstlisting}
random_integers_between_150_and_300 = random_integers_between_50_and_100 * 3
print(random_integers_between_150_and_300)
\end{lstlisting}

%==============================================================================
\section{Bài tập 1: Tạo dữ liệu tuyến tính}
%==============================================================================

\begin{phuongphap}[title={Yêu cầu bài tập}]
Tạo dataset đơn giản gồm một đặc trưng (feature) và một nhãn (label):
\begin{enumerate}
  \item Gán dãy số nguyên từ 6 đến 20 (bao gồm cả 6 và 20) vào NumPy array tên \texttt{feature}.
  \item Gán 15 giá trị vào NumPy array tên \texttt{label} sao cho:
  \[
    \texttt{label} = 3 \cdot \texttt{feature} + 4.
  \]
\end{enumerate}
Ví dụ phần tử đầu tiên của \texttt{label} là:
\[
  3 \cdot 6 + 4 = 22.
  \]
\end{phuongphap}

\begin{lstlisting}
feature = ? # write your code here
print(feature)
label = ?   # write your code here
print(label)
\end{lstlisting}

\begin{lstlisting}
feature = np.arange(6, 21)
print(feature)
label = (feature * 3) + 4
print(label)
\end{lstlisting}

%==============================================================================
\section{Bài tập 2: Thêm nhiễu (noise) vào dataset}
%==============================================================================

\begin{phuongphap}[title={Yêu cầu bài tập}]
Để dữ liệu thực tế hơn, hãy thêm một chút nhiễu ngẫu nhiên vào từng phần tử của \texttt{label}:
\begin{itemize}
  \item Mỗi phần tử của \texttt{label} được cộng với một số thực ngẫu nhiên \textbf{khác nhau} trong đoạn \([-2, 2]\).
  \item Không dựa vào broadcasting. Hãy tạo một array \texttt{noise} có cùng kích thước với \texttt{label}.
\end{itemize}
\end{phuongphap}

\begin{lstlisting}
noise = ?    # write your code here
print(noise)
label = ?    # write your code here
print(label)
\end{lstlisting}

\begin{lstlisting}
noise = (np.random.random([15]) * 4) - 2
print(noise)
label = label + noise
print(label)
\end{lstlisting}

\end{document}

